% ============================================================
%  Back cover — Whatever We May Be
%  Cédric Laubscher — 2026
% ============================================================
\documentclass[12pt,a4paper]{article}
\usepackage[utf8]{inputenc}
\usepackage[T1]{fontenc}
\usepackage[british]{babel}
\usepackage{csquotes}
\usepackage[sfdefault]{sourcesanspro}
\usepackage[top=3.5cm,bottom=3.5cm,left=3cm,right=3cm]{geometry}
\usepackage{amsmath}
\setlength{\parindent}{0pt}
\setlength{\parskip}{0.9em}
\pagestyle{empty}
\begin{document}

\begin{center}
{\large\itshape Whatever We May Be}\\[0.3em]
{\small An Introduction to the Projective Dynamic Logo}
\end{center}

\bigskip

Why does an electron exist? Not how it behaves---physics answers that with extraordinary precision. But \emph{why} that particular form, that absolute stability, that perfect identity across all of space and time? This \emph{why} is the crack that physics has not managed to close in a century.

The Projective Dynamic Logo suspends everything: space, time, particles. And it poses the problem in its absolute nakedness: under these conditions, what can exist? The answer is simple and vertiginous. For something to exist durably, it is necessary and sufficient that it form a cycle---a difference that returns to itself. A pulse. From this single principle, without presupposing either space or time, the PDL reconstructs the logical form of the electron, then of the proton, then the great constants of physics---$\hbar$, $\alpha$, $k_{\rm B}$, $G$---no longer as numbers fallen from the sky, but as signatures of the way each structure manages its own internal equilibrium.

But no complex structure can close perfectly on itself. There always remains an incompressible leakage---an irreducible opening towards what lies beyond. It is from this leakage that gravitation arises. It is this same opening that makes freedom possible, transcendence possible---and our vulnerability.

An electron, a proton, a human being: the necessary instruments of one and the same implacable logic. This book offers the first unified reading, from the axioms to the human condition.

\bigskip

\begin{center}
{\small\itshape All technical publications of the PDL are freely available on \textbf{Zenodo}, under the name Cédric Laubscher.}
\end{center}

\end{document}